\subsection{A conditional completion of the entire coefficient-five third gate}

\begin{result}{E.205}{Proved}
{The first finite five-primary system suffices for a controlled original continuation}
Fix a legal first pair $q_0,r_0$ in the minimal-five-shell
subfamily E.199 of the unchanged original recursion E.102.
Suppose the finite affine system E.202.1 is consistent for this pair.
Then there is a continuation of that original recursion for which
the ordinary infinite-support pair has an integral class-three lift
in the geometric pro-five quotient.  Equivalently, the coefficient-five
component of the continuous third obstruction of its integral
class-two section is zero.

More precisely, retain $n_j=(j+4)!$ and put
\[
 c_j=v_5(n_j),\qquad t_j=v_5(j+5),\qquad
 e_j=v_5(n_{j+1})=c_j+t_j.
 \tag{E.205.1}
\]
The continuation can be chosen with the E.199 shells at every stage
and, in the literal coordinates of B.102, with
\[
 \alpha_j\in\mathbb Z_5^\times,\quad
 \beta_j\in\mathbb Z_5^\times,\quad
 D_{A,j},D_{B,j}\in5^{c_j}\mathbb Z_5
                         \qquad(j\geq1).
 \tag{E.205.2}
\]
Here $\alpha_j$ is the first abelian Frobenius coordinate at $q_j$
and $\beta_j$ is the second at $r_j$.
Consequently the two normalized sequences satisfy
\[
 v_5(T_{A,j}),\ v_5(T_{B,j})\ \geq\ c_j-t_j
                  \longrightarrow+\infty\qquad(j\longrightarrow\infty).
 \tag{E.205.3}
\]
For a fixed legal first pair, existence of an original continuation
passing the full continuous coefficient-five third gate is therefore
equivalent to consistency of E.202.1. E.213 subsequently proves
that such legal first pairs exist, and E.211 strengthens this
continuation to class-two support matching. These results leave
the other coefficient primes and higher pro-five matching
unresolved; they do not establish an all-depth section.
\end{result}

\begin{proof}
We prove the extra prime choices, the finite cochain stability, and
the delayed correction separately.  All local roots below remain
roots of the original exponent-$n_j$ equation; no factorial or
complete-splitting requirement is changed.

\emph{Two independent nonzero Kummer conditions.}
Put
\[
                  A_0=3q_0^{24},\qquad B_0=-2r_0^{24}.
 \tag{E.205.4}
\]
For every continuation, its literal abelian cocycles modulo five
are the Kummer cocycles of these two numbers.  At every stage
$j\geq1$ we may require, in addition to all the E.102 conditions,
\[
 A_0\text{ is not a fifth power modulo }q_j,\qquad
 B_0\text{ is not a fifth power modulo }r_j.
 \tag{E.205.5}
\]
Here is a disjointness proof over the full cyclotomic field used
by that stage, so these requirements are not an appeal to
generic prime choices.

At the first prime choice write $a=a_j,b=b_j,n=n_j$ and
$m=n_{j+1}$ as in E.102.  Let $F$ be the cyclotomic field
containing every CRT modulus, $\mu_m$, and every other
cyclotomic mark required at this step.  The prescribed CRT
automorphism of $F$ is the identity on $\mu_m$.  Let
\[
                       R=F(b^{1/n}).
\]
The original radical argument of E.102, applied after enlarging
the cyclotomic field, gives a cyclic group
$\operatorname{Gal}(R/F)$ of order $n/2$.
In particular its unique subextension of degree five is
$F(b^{1/5})=F(B_0^{1/5})$; all later exponents in $b$ are
divisible by five.  This degree-five subextension is unramified
above three.  Indeed $b$ is a unit there and $X^5-b$ has unit
derivative at every root in residue characteristic three, so its
root algebra is finite etale over the local valuation ring.

On the other hand $v_3(A_0)=1$.  The ramification index at three
of any cyclotomic field over $\mathbb Q$ divides
$2\cdot3^a$ for some $a$, and is prime to five.  Thus the
valuation of $A_0$ in $F$ is not divisible by five, and
$F(A_0^{1/5})/F$ is totally ramified of degree five there.
It cannot be the indicated subextension of $R$, whence
\[
                       F(A_0^{1/5})\cap R=F.
 \tag{E.205.6}
\]
The E.102 Frobenius on $R$, with its prescribed restriction to
$F$ and its fixed chosen $n$th root of $b$, can therefore be
combined with any of the five Frobenius translations on
$A_0^{1/5}$.  Choose a nonzero one.  These automorphisms agree
on $F$, and the CRT automorphism fixes $\mu_5$; the nonzero
translation condition is invariant under conjugacy.  Chebotarev
in precisely the form used in E.102 supplies the required $q_j$,
still beyond every original size bound and exclusion.

For the second prime choice put $a'=a_jq_j^{n_j}$ and let $F$
now include the enlarged CRT modulus, including the specified
residue at $q_j$.  The prescribed radical is
$F((a')^{1/m})$, whose group over $F$ is cyclic of order $m/2$.
Its unique degree-five subextension is
$F((a')^{1/5})=F(A_0^{1/5})$, unramified above two since
$a'$ is a two-adic unit.  In contrast, $v_2(B_0)=1$ and a
cyclotomic ramification index at two is a power of two.
Thus $F(B_0^{1/5})/F$ is totally ramified of degree five above
two and is disjoint from that radical extension over $F$.
The same gluing and Chebotarev argument imposes the second
condition in E.205.5 while preserving the original complete
splitting and every current-prime congruence.

For clarity, the degree assertion just used does not require
the enlarged cyclotomic field to be unramified at any other
prime.  A radical of the forms $3w^8$ or $-2w^8$ has no
odd prime root or fourth root in any cyclotomic field: the
private valuation and the nonnormal pure-root subfield argument
in E.102 prove this.  It is already a square there after
adjoining $\mu_{24}$.  Kummer theory therefore gives exactly
the degree $n/2$ or $m/2$ over the enlarged $F$, as asserted.
The sole degree-five intersection is then excluded by the
explicit local ramification comparison.

Since $5\mid n_j$ for $j\geq1$, the value of an abelian
Frobenius modulo five at the newly ramified prime is independent
of the choice of its tame lift: changing that lift adds a
multiple of $n_j$.  Every tail abelian term is divisible by
five.  Thus E.205.5 gives exactly the two unit assertions in
E.205.2 for the eventual integral cochains.

\emph{A finite cochain invariant after each actual prefix.}
Fix $j\geq1$, after the actual primes through $q_j,r_j$ have
been chosen.  Every earlier support pair $q_i,r_i$, $i<j$, has
already received its one delayed root correction.  Require both
old roots at these earlier places to be congruent to one modulo
that place at every later stage.  This requirement is legal:
the preceding correction stage made the coefficient sum one
modulo the old prime; simple-root lifting applies until that
prime divides the factorial, and the stronger original Newton
estimate applies thereafter.  This is exactly the retention
argument in E.202.  It imposes conditions at a finite set at
each stage and does not affect the separate choices at five.

Consider any two complete legal continuations with this prefix
and these retention conditions.  At $q_j,r_j$ allow arbitrary
old roots at the next stage; after that stage require the same
principal-unit retention there.  Such continuations exist by
E.102 and the just-proved additional prime-selection argument.
Choose the global integral Kummer cochains as the uniformly
convergent sums with the displayed factorial coefficients.
Their reductions modulo $5^{c_j}$ agree literally on the two
continuations.  If $h,h'$ are their integral class-two
cochains, then
\[
                  z=h'-h\pmod{5^{c_j}}
 \tag{E.205.7}
\]
is a cocycle with coefficients $\mathbb Z/5^{c_j}(2)$.
The existence of these integral class-two sections follows
from A.25, independently of the third obstruction.

We check every possible tame residue of $z$ outside five.
At an earlier support prime the ratio of each later selected
prime to its comparator is a principal unit.  Its local
five-primary Kummer cochain has an integral primitive.  A
factor carrying exponent $n_k$, $k\geq j+1$, therefore
changes the complementary divided primitive by a multiple
of $n_k$ in $\mathbb Z_5$; the convergent sum of all such
changes is divisible by $5^{c_j}$.  The exact integral
class-two tame formulas
\[
            h(T_{q_i})=n_i u_i,\qquad
            h(T_{r_i})=-n_i v_i
 \tag{E.205.8}
\]
make the corresponding changes of central inertia zero
modulo $5^{c_j}$, even when $n_i$ itself has smaller
five-adic valuation.

At either current support prime $q_j,r_j$, and at every
new support prime of either continuation, the ordinary
valuation exponent is divisible by $5^{c_j}$.  The same
formulas give zero central inertia modulo that power,
regardless of the complementary root choice.  At a place
which is support in only one continuation the other
central inertia is zero as well.  At zero-support places
the integral class-two tame relation forces central
inertia zero.  Finally at two and three the primitive
comparison uses division by $p-1$, which is a five-adic
unit; the identical abelian cochains modulo $5^{c_j}$
therefore give identical central inertia modulo that
power.  This accounts for every finite place outside five.

Thus $z$ is unramified outside $\{5,\infty\}$.  The
weight-two groups in E.202.2 and the multiplication-by-$5^{c_j}$
coefficient sequence give
\[
 H^1(G_{\mathbb Q,\{5,\infty\}},\mathbb Z/5^{c_j}(2))=0.
 \tag{E.205.9}
\]
Consequently $z$ is a central coboundary.  Its Frobenius
value at a current support prime $p=q_j$ or $r_j$ is
$(p^2-1)$ times a central scalar, and vanishes modulo
$5^{c_j}$ because the minimal shells give
$v_5(p^2-1)=e_j\geq c_j$.
It follows that
\[
       \gamma_j\pmod{5^{c_j}},\qquad
       \widetilde\gamma_j\pmod{5^{c_j}}
 \tag{E.205.10}
\]
are already determined by the actual prefix and the
retention conditions.  No final infinite cochain is used
to choose that prefix.  One may use any legal completion
to define these two invariant values; the preceding
argument proves independence of the choice.  The argument
only makes $h$ fixed up to a coboundary, which is exactly
what E.205.10 requires.

\emph{The one-stage delayed correction.}
At stage $j+1$, the old local equation at $q_j,r_j$ has
exponent $n_{j+1}$.  Both local fields contain
$\mu_{5^{e_j}}$, by their minimal shells.  Choose any old
unit-coordinate solution.  At $q_j$ multiply its second
coordinate by $\zeta_{5^{e_j}}^i$, and at $r_j$
independently multiply its first coordinate by
$\zeta_{5^{e_j}}^k$.  These are still exact roots of the
same equation.  The original size bounds give
$q_j,r_j>n_{j+2}$, so neither old prime divides the next
factorial $n_{j+2}$; no original condition requiring these
roots to be one modulo a divisor of that factorial is
violated.  All other old roots, including the five-adic
shell roots, are retained.

In compatible local Tate coordinates, the Kummer primitive
of $\zeta_{5^{e_j}}$ is $1/5^{e_j}$ modulo
$\mathbb Z_5$.  The selected next-prime residue has the
specified root-of-unity factor, with only a principal-unit
ambiguity.  Multiplication by its actual exponent
$n_{j+1}$ therefore changes the divided primitives by
\[
 \Delta u_j=\frac{n_{j+1}}{5^{e_j}}i
                  \pmod{5^{c_j}},\qquad
 \Delta v_j=\frac{n_{j+1}}{5^{e_j}}k
                  \pmod{5^{c_j}}.
 \tag{E.205.11}
\]
The coefficient is a five-adic unit, and $e_j\geq c_j$,
so each primitive can be prescribed arbitrarily modulo
$5^{c_j}$.  The two choices concern distinct old local
fields and are independent CRT data.  The prime-selection
argument above still supplies the next primes with
E.205.5 after these root choices.

By E.205.10 the two central Frobenius values are unchanged
modulo $5^{c_j}$, whereas the two abelian coefficients
$\alpha_j,\beta_j$ are units.  We can therefore choose
$i,k$ so that
\[
 \gamma_j-q_j\alpha_j u_j=0\pmod{5^{c_j}},\qquad
 \widetilde\gamma_j+\beta_jv_j=0\pmod{5^{c_j}}.
 \tag{E.205.12}
\]
After this stage both coefficient sums are one modulo
their old support prime.  Require principal-unit roots
there at every subsequent stage.  The same primitive
comparison and coboundary argument show that E.205.12
then survives all later choices, including the later
root-of-unity corrections at other support primes.

The initial pair is handled separately.  Consistency of
E.202.1 selects one of its twenty-five legal stage-one
root choices, after which its two mod-five conditions
are retained by principal-unit roots.  The additional
nonzero Kummer conditions on $q_1,r_1$ are compatible
with each of those root choices, as proved above.
Thereafter stage $j+1$ corrects the pair selected at
stage $j$, for each $j\geq1$.  At each finite stage
only finitely many old-place conditions are imposed.
The original induction, local unit-equation invariant,
and all completions therefore remain valid, giving an
ordinary infinite-support pair with E.205.2.

\emph{The strict global condition.}
For $j\geq1$ one has $c_j\geq t_j$ and
$c_j-t_j\longrightarrow\infty$.  To check the early
indices explicitly, for $j=1,2,3,4$ the triple
$(c_j,t_j,e_j)$ is $(1,0,1)$; for $j=5$ it is
$(1,1,2)$, and for $j=20$, the first $j+5$ divisible
by twenty-five, it is $(4,2,6)$.  Generally, writing
$M=j+5\geq6$, Legendre's factorial formula gives
\[
 c_j=\sum_{a\geq1}\left\lfloor\frac{M-1}{5^a}\right\rfloor,
 \qquad t_j=v_5(M).
\]
If $t_j=0$ there is nothing to prove.  If $t_j=1$
then $M\geq10$ and $c_j\geq1$.  If $t_j\geq2$
then $c_j\geq5^{t_j-1}-1\geq t_j$.
Finally $c_j\geq(M-1)/5-1$ and
$t_j\leq\log_5M$, proving divergence of the difference.

The units $d_{A,j},d_{B,j}$ of E.199 give exactly
\[
 T_{A,j}=\frac{D_{A,j}}{(j+5)d_{A,j}},\qquad
 T_{B,j}=\frac{D_{B,j}}{(j+5)d_{B,j}}.
\]
Equations E.205.2 therefore imply both integrality and
the zero limits E.205.3.  The first two terms are integral
by the chosen E.202 parameter.  At two and three the
coefficient-five local third groups vanish; the group
at five also vanishes because the mod-five cyclotomic
character squared is nontrivial.  The real group has
no five-primary positive cohomology.  Zero-support tame
residues vanish by B.102.  These are exactly the remaining
local places already accounted for in E.199.

The extra normalized polynomial terms of B.102 are
integral and tend to zero, including the actual action
term $n_j/24$.  The strict continuous localization
comparison A.22 and B.22--B.23 now makes the entire
coefficient-five third obstruction zero, rather than
merely giving zero at every local place.  Here there is
no remaining fixed-ramification kernel:
$H^2(G_{\mathbb Q,\{5,\infty\}},\mathbb Z_5(3))=0$
by the weight-three $K_4$ calculation in A.22.
Standard
continuous central-extension obstruction theory gives
the stated global integral pro-five class-three lift.
Conversely any original continuation with that lift
passes the first two local tests, and E.202's exhaustion
forces consistency of its finite system.  This proves
the equivalence and its stated scope.
The exact factorial indices and finite affine correction
maps are additionally checked by
\path{code/ghost_audit_five_continuation.py}; its saved
output explicitly makes no legal-first-pair assertion.
\end{proof}
